\documentclass[11pt,letterpaper]{article}
\usepackage[letterpaper,left=.62in,right=.62in,top=.64in,bottom=.62in,headheight=14pt,headsep=13pt,footskip=22pt]{geometry}
\usepackage[T1]{fontenc}
\usepackage{lmodern}
\usepackage{amsmath,amssymb,array,tabularx,fancyhdr,url}
\pdfmapfile{+lm.map}
\pdfcompresslevel=9
\pdfobjcompresslevel=3
\renewcommand{\familydefault}{\sfdefault}
\pagestyle{fancy}\fancyhf{}
\fancyhead[L]{\fontsize{10}{12}\selectfont Week 78 / Three-armed lines / Adult guide}
\fancyfoot[L]{\fontsize{9}{11}\selectfont Bellingham Math Circle / Week 78 / F78-FAC-v1}
\fancyfoot[R]{\fontsize{9}{11}\selectfont\thepage}
\renewcommand{\headrulewidth}{0pt}\renewcommand{\footrulewidth}{0pt}
\setlength{\parindent}{0pt}\setlength{\parskip}{4pt}
\setlength{\tabcolsep}{5pt}\renewcommand{\arraystretch}{1.18}
\AtBeginDocument{\fontsize{10.5}{12.8}\selectfont}
\newcommand{\head}[1]{\par\vspace{4pt}{\large\bfseries #1}\par\vspace{1pt}}
\newcommand{\prob}[2]{\par\vspace{4pt}\textbf{Problem #1. #2}\par}
\newcommand{\hints}[1]{\par\textbf{Hints, in order.} #1\par}
\newcommand{\LJ}{L(a,b)}
\begin{document}
\head{The mathematical destination}
\textbf{Intersection theorem.} Every two lines in this packet meet. With distinct junctions they share exactly one point, except for three alignments, when they share a whole ray. Equal junctions give identical lines. Here a line $L(a,b)$ means the three closed rays
\[
 (a+t,b),\quad (a,b+t),\quad (a-t,b-t),\qquad t\geq0.
\]
Coordinates may be any real numbers; $t$ is any nonnegative real number. East, north and southwest stay fixed; moving a line means translating it without turning it. The rays continue beyond the paper.

\par\noindent\begin{tabularx}{\linewidth}{@{}p{2.43in}X@{}}
\textbf{Junctions $A=(a,b)$, $B=(c,d)$} & \textbf{Their complete intersection}\\\hline
Same junction & The whole three-armed line\\
Same column, different junctions & North ray from the higher junction\\
Same row, different junctions & East ray from the rightmost junction\\
Same rising diagonal: $a-b=c-d$, different junctions & Southwest ray from the more southwesterly junction\\
None of these & Exactly one point
\end{tabularx}

\textbf{Joining theorem.} Two \emph{different} target points lie on exactly one such line unless they share a row, column or rising slope-1 diagonal. In each aligned case, infinitely many lines work. All possible junctions form, respectively, a west ray from the left target, a south ray from the lower target, or a northeast ray from the northeast target.

\textbf{Why this is tropical geometry.} $L(a,b)$ is where the least of $x-a$, $y-b$, and $0$ occurs at least twice. It is the planar crease set of their minimum. A tie between two larger numbers does not count. Adding the same constant to all three scores preserves the set. The opening example uses $\min(x,y,2)$ and junction $(2,2)$.

\textbf{Discovery and limits.} Problems 1--2 distinguish a crossing from an overlap; 4--5 reverse the task to find a junction. Problems 3, 6 and 7 ask for arguments covering all placements: offer the proof on page 4 when needed. Finite trials can suggest a rule, but do not prove it. We use ordinary set intersection throughout; a shared ray contains infinitely many points. Stable intersection is a different notion and is outside this session.

\head{Target table and preparation}
\textbf{Scope.} Three Grades 4--5 children with the organizer. Children need to plot points on a square grid, compare small numbers, and retain the fixed directions. Signed arithmetic is unnecessary for entry. A useful first session ends with correct drawings and an explained crossing/overlap distinction. General proof may take another visit. The other eight children and their two adults need separately prepared work; this prototype supplies no K--1 or Grades 2--3 route.

\textbf{Make two working pairs:} two children together, the third with the adult. Rotate the adult's partner after each board. The adult reads or steadies a ruler as needed; children choose placements and draw.

\textbf{Per pair:} one shared student packet, two contrasting coloured pencils, one ruler, two junction counters, one eraser, two large square-grid sheets. \textbf{Table total including spares:} three student packets (12 single-sided sheets); six coloured pencils (two spare); three rulers; six counters; three erasers; six grid sheets. Print one copy of this four-page guide. Print student pages 1--4 at 100\% on Letter, but hand out one page per pair at a time; keep the third packet as a spare.

\textbf{Beforehand:} on each extra grid sheet mark $0$--$12$ on both axes, using 1\,cm squares. No cutting or taping is required. Optional: two tracing sheets per pair plus two spare (six total), for drawing and translating the same three-arm shape. Test pencil contrast, ruler access, grid scale and visibility through tracing paper with the actual supplies. Rehearse drawing a southwest ray and continuing off the small printed grid. \textbf{Physical rehearsal and classroom piloting have not been performed.}
\newpage
\head{The hour}
\par\noindent\begin{tabularx}{\linewidth}{@{}p{.72in}X@{}}
0--5 min & Running around, then settle with the materials.\\
5--8 & Brief whole-group demonstration below; the other tables then use their separately prepared activities.\\
8--23 & Problem 1: choose, draw, and compare. Rotate partners after a board.\\
23--36 & Problem 2: see all three kinds of shared ray.\\
36--43 & Problem 3 orally or in a drawing. A conjecture is a valid stopping point if labelled as such.\\
43--53 & If ready, Problem 4 then 5; otherwise the fallback below. Problems 6--7 are a continuation for a quick pair or another session.\\
53--60 & Share one example, record a claim, and put materials away.
\end{tabularx}

\textbf{Launch: exactly one legal action.} Put a counter at $(2,2)$ on a large square grid. Say, ``From this junction I draw right, up, and diagonally down-left. Each arm goes on forever.'' Draw all three with arrows. Point to $(4,2)$: ``The numbers are $4,2,2$; the smallest is tied.'' Point to $(4,4)$: ``Here they are $4,4,2$; the smallest is alone.'' Ask one child to place the counter at $(3,2)$ and draw the same three directions. The others check that the shape did not turn. Touch $(5,2)$ on its east arm. Do not announce the intersection theorem.

\textbf{Fallback at this table.} One child chooses two junctions on a large grid; the partner draws one line while the chooser draws the other. The third child checks directions and marks what is shared. Swap roles and try to make a different kind of meeting. Drop coordinate recording if it gets in the way. \textbf{Closing question:} ``Which drawing changed your mind about how these lines can meet?''

\head{Student page 1: choice and crossings}
\prob{1}{Core. Two chosen pairs, then two fixed pairs.}
Children choose $B$ on the two top grids, where $A=(4,4)$. Many answers work. Ask for two different meeting types; do not select $B$ for them. This complete menu helps check their choices:

\par\noindent\begin{tabularx}{\linewidth}{@{}p{1.35in}X@{}}
\textbf{Example choice of $B$} & \textbf{Complete common part with $L(4,4)$}\\\hline
$(6,5)$ & Single point $(5,4)$\\
$(6,4)$; $(4,6)$ & East ray from $(6,4)$; north ray from $(4,6)$\\
$(6,6)$ & Southwest ray from $(4,4)$\\
$(4,4)$ & The whole line; the problem permits equal junctions
\end{tabularx}

\textbf{Fixed bottom grids:} left, $(2,2)$ and $(6,5)$ meet only at $(3,2)$; right, $(2,2)$ and $(5,6)$ meet only at $(2,3)$. In each, the upper junction's southwest arm meets an east or north arm of the lower line.
\hints{1. ``Trace the three arms from each counter. Did either shape turn?'' 2. Follow one coloured arm and look for the other colour on it. 3. If only crossings appear, ask what happens when $B$ stays on $A$'s row.}
\textbf{If struggling:} do one chosen grid and one bottom grid, then use the large board for Problem 2. Preserve the child's choice of $B$.

\head{Student page 2: overlapping rays}
\prob{2}{Core. Mark the whole common part.}
Reading left to right, top to bottom:
\par\noindent\begin{tabularx}{\linewidth}{@{}p{2.0in}X@{}}
$(2,5)$ and $(6,2)$ & One point, $(6,5)$\\
$(3,2)$ and $(3,6)$ & Infinitely many: $(3,y)$ for every $y\geq6$\\
$(2,4)$ and $(6,4)$ & Infinitely many: $(x,4)$ for every $x\geq6$\\
$(2,2)$ and $(6,6)$ & Infinitely many: $(2-t,2-t)$ for every $t\geq0$
\end{tabularx}
\hints{1. ``Where can you trace over both colours at once?'' 2. Continue that trace past the page edge. 3. Place a finger between two grid dots: does that point belong to both too?}
\textbf{If struggling:} compare the top two grids first. Say ``a whole ray, endlessly many points''; do not count visible grid dots as the answer.
\newpage
\head{Solutions and hints, continued}
\prob{3}{Core discussion; full proof can wait.}
\textbf{Answer: no.} With different junctions the full common set is either one point or a ray. It cannot be exactly two points. Child explanation after the rectangle argument: ``If the junctions are on the same row, column or diagonal, the matching arms keep sharing forever. Otherwise the arms meet once at the rectangle crossing. Those cover every placement.'' Do not accept ``we could not draw one'' as an impossibility proof.
\hints{1. ``Can you point to exactly two shared points without sharing the stretch between them?'' 2. Put a rectangle around the two junctions. 3. Work through its tall, wide and square forms using page 4.}
\textbf{Stopping point:} save a child's proposed rule even if its explanation is unfinished.

\prob{4}{Further. Construct a line through two targets.}
\textbf{Left:} $P=(2,5)$, $Q=(6,2)$ force junction $(2,2)$. $P$ is on the north arm and $Q$ on the east arm. \textbf{Right:} $P=(2,2)$, $Q=(6,5)$ force junction $(5,5)$. $P$ is on the southwest arm and $Q$ on the east arm. Neither pair permits a second line, even with a noninteger or off-page junction. The two reversed allowed-junction figures meet only at that junction (page 4).
\hints{1. ``Move a counter for the junction. Can one unchanged shape touch both targets?'' 2. Hold one target fixed and try putting it on each of the three arms. 3. Mark every junction allowed by that target, then do the same for the other target.}
\textbf{If struggling:} construct the left example only; let the child check both arms instead of asking for uniqueness now.

\prob{5}{Further. Find all junctions, including between dots.}
Reading left to right, these are the \emph{complete} sets:
\par\noindent\begin{tabularx}{\linewidth}{@{}p{1.76in}X@{}}
$P=(2,4)$, $Q=(6,4)$ & $(2-t,4)$, $t\geq0$: every point west of $(2,4)$, including $(2,4)$.\\
$P=(3,2)$, $Q=(3,6)$ & $(3,2-t)$, $t\geq0$: every point south of $(3,2)$, including $(3,2)$.\\
$P=(2,2)$, $Q=(5,5)$ & $(5+t,5+t)$, $t\geq0$: every point northeast of $(5,5)$, including $(5,5)$.
\end{tabularx}
All three sets are infinite. The targets share an east, north or southwest arm, respectively. Sliding along the stated ray keeps both covered. The reversed figures on page 4 have no other common junctions, proving completeness. Include the endpoint, points between dots and continuation beyond the page.
\hints{1. ``Can you slide a working junction a little while both targets stay on the line?'' 2. Slide the other way until it stops working. 3. Mark the allowed-junction figure for each target to see where both requirements hold.}
\textbf{If struggling:} keep only the same-row example; finding its full west ray is enough.

\prob{6}{Continuation. Off-page intersection and the general rule.}
\textbf{Answer:} $A=(2,7)$ and $B=(10,5)$ meet only at $(10,7)$: east from $A$, north from $B$. It is outside the printed window. No two lines miss. More than one shared point occurs exactly for a shared column, row or rising diagonal; for equal junctions the whole lines coincide. The rectangle proof (page 4) uses directions and lengths, covering noninteger and off-page junctions too.
\hints{1. ``Continue the drawing on a large grid. Where is $B$?'' 2. Draw the rectangle with $A,B$ as opposite corners. 3. Compare different rectangle shapes, including zero width, zero height and equal width and height.}
\textbf{If time is short:} find the off-page point and save the general proof.

\prob{7}{Continuation. How many lines through two different points?}
\textbf{Answer:} exactly one unless the points share a row, column or rising slope-1 diagonal; in each exception, infinitely many. Page 1 gives all exceptional junction sets. The reversed-junction proof (page 4) rules out zero lines or a finite number greater than one.
\hints{1. ``Keep one target fixed. Where can the junction move?'' 2. Trace its three allowed directions, then repeat for the other target. 3. Turn both \emph{allowed-junction drawings} through a half-turn and use the proved intersection theorem.}
\textbf{If not ready:} retain the contrasting examples from Problems 4--5 and revisit later.
\newpage
\head{The mathematics behind the answers}
\textbf{The three-arm shape.} Put $u=x-a$, $v=y-b$. A tied minimum of $u,v,0$ means one of
\[
 u=0\leq v\quad(\text{north}),\qquad
 v=0\leq u\quad(\text{east}),\qquad
 u=v\leq0\quad(\text{southwest}).
\]
These cases exhaust every tie; all three meet at the junction. This proves the shape and its orientation. Subtracting $2$ from $x,y,2$ gives the opening junction $(2,2)$ without changing which scores are smallest. A common added constant preserves every comparison.

\textbf{Why two lines meet, and why there is no extra meeting.} Draw the closed axis-parallel rectangle between their junctions. It may have zero width or height. Exchange the names of the junctions when needed.

\textbf{Northwest/southeast corners.} The northwest junction's east arm meets the southeast junction's north arm at the northeast corner of the rectangle. Their other north/east arms point away from each other. Each southwest arm goes below or left of the other junction's relevant north/east arm; the two southwest arms are distinct parallels. Thus there is exactly one meeting.

\textbf{Southwest/northeast corners, with positive width and height.} Starting at the northeast junction, follow its southwest arm. If the rectangle is taller than wide, the arm first hits its left edge, on the other line's north arm. If wider than tall, it first hits its bottom edge, on the other line's east arm. After exiting that side it cannot reach the other side-ray in its allowed direction. The northeast junction's east and north arms cannot reach the southwest line. The southwest arms are parallel and distinct. Thus each unequal-sided rectangle gives just one meeting.

\textbf{The square and the boundary cases.} For a square, the northeast junction's southwest arm reaches the southwest junction and continues along its entire southwest ray. No other arms add common points. For zero width, the north rays coincide above the higher junction; the remaining arms have no extra meetings. For zero height, the east rays coincide to the right of the rightmost junction, with no extra meetings. If both dimensions vanish, the two whole lines coincide. These cases exhaust all real placements and prove the table on page 1. In particular, two distinct lines cannot share exactly two points or just a bounded nonzero segment.

\textbf{Why targets determine the stated junctions.} Fix a target $P=(r,s)$. If it is on an east arm, the junction is $(r-t,s)$; if on a north arm, $(r,s-t)$; if on a southwest arm, $(r+t,s+t)$, with $t\geq0$. Conversely, each such junction works. Thus the full allowed-junction set has west, south and northeast rays. This reversed figure is a tool for locating junctions, not one of our original min-plus lines.

A line through both targets corresponds exactly to a common point of their two allowed-junction figures. Reflect both figures through the origin. They become the original fixed-orientation family, so the intersection theorem applies. Reflect back: different unaligned targets have one junction; aligned targets have precisely the west, south or northeast ray stated on page 1. Each different junction defines a different line, proving the joining theorem. Repeated targets, excluded in Problem 7, would allow the whole reversed three-arm figure of junctions.

\head{After the session}
Record which problems and boards each pair reached; a placement or drawing worth keeping; the child's actual words; and any adult hint that changed the work. Label each claim \textbf{tried}, \textbf{guessed}, \textbf{checked in some cases}, or \textbf{explained for all placements}. Note whether fixed directions, ray continuation, between-dot points or the inverse task caused trouble. Do not infer learning from a completed drawing alone. Record material changes and any rehearsal separately from classroom evidence.

\head{Source lineage and next questions}
{\fontsize{9.5}{11.6}\selectfont
Speyer and Sturmfels, \emph{Tropical Mathematics} (2004), \S1 and \S3, supplies the min-plus arithmetic and tropical-line framework:\par
\url{https://www.claymath.org/wp-content/uploads/2022/03/Speyer.pdf}\par
Tewari, \emph{Point-Line Geometry in the Tropical Plane} (2020), \S3--4, is further reading on incidence and the distinction between ordinary and stable intersection. It uses max-plus, whose arms are reversed from this packet:\par
\url{https://arxiv.org/abs/2006.04425}\par
The packet's numerical cases, rectangle proof and practical route are independently developed. The guide checker solves every pair of parameterized arms using exact fractions, separately from the student checker; its finite audit supports the examples, not the all-real theorem. A later question is which triples of targets can lie on one line: begin with a pair having a unique junction and test the third target. The reversed-junction construction is an elementary instance of point-line duality. No classroom outcome is claimed.
}
\end{document}
